\documentclass[12pt]{article}
% Shared preamble and text for the data-request letters.
% Replace the bracketed fields before submitting any letter.

\usepackage[a4paper,top=2.4cm,bottom=2.4cm,left=2.8cm,right=2.4cm]{geometry}
\usepackage[utf8]{inputenc}
\usepackage[T1]{fontenc}
\usepackage{lmodern}
\usepackage{microtype}
\usepackage{setspace}
\usepackage{array}
\usepackage{booktabs}
\usepackage{enumitem}
\usepackage{tabularx}
\usepackage{xcolor}
\usepackage[hidelinks]{hyperref}

\setstretch{1.12}
\setlength{\parindent}{0pt}
\setlength{\parskip}{6pt}
\setlist[itemize]{leftmargin=1.5em,itemsep=2pt,topsep=2pt}
\setlist[enumerate]{leftmargin=1.7em,itemsep=2pt,topsep=2pt}
\newcolumntype{Y}{>{\raggedright\arraybackslash}X}
\newcolumntype{P}[1]{>{\raggedright\arraybackslash}p{#1}}

\newcommand{\NamaMahasiswa}{[NAMA LENGKAP MAHASISWA]}
\newcommand{\NIMMahasiswa}{[NIM]}
\newcommand{\NamaPembimbing}{[NAMA DAN GELAR DOSEN PEMBIMBING]}
\newcommand{\JudulSementara}{Diagnosis \textit{Borrowed Size} dan \textit{Agglomeration Shadow} antara Jakarta dan Pusat-pusat Sekunder di Kawasan Aglomerasi Jakarta}

\newcommand{\KopSurat}{%
  \begin{center}
    \fbox{%
      \begin{minipage}{0.88\textwidth}
        \centering
        \textbf{KOP SURAT RESMI FAKULTAS TEKNIK UNIVERSITAS GADJAH MADA}\\
        \small Ganti blok ini dengan kop resmi sebelum surat dikirim.
      \end{minipage}}
  \end{center}
  \vspace{4pt}
}

\newcommand{\IdentitasMahasiswa}{%
  \begin{tabularx}{\textwidth}{@{}P{0.30\textwidth}Y@{}}
    Nama & \NamaMahasiswa \\
    NIM & \NIMMahasiswa \\
    Program studi & Sarjana Perencanaan Wilayah dan Kota \\
    Departemen dan fakultas & Departemen Teknik Arsitektur dan Perencanaan, Fakultas Teknik, Universitas Gadjah Mada \\
    Dosen pembimbing & \NamaPembimbing
  \end{tabularx}
}

\newcommand{\PernyataanPenggunaanData}{%
Data yang dimohon akan digunakan untuk kepentingan pendidikan dan penelitian nonkomersial. Data tidak akan dijual, dialihkan, atau disebarluaskan dalam bentuk mentah. Hasil yang dipublikasikan akan berupa statistik agregat, peta turunan, atau keluaran lain yang diizinkan oleh pemilik data. Pengolahan dan penyimpanan data akan mengikuti ketentuan kerahasiaan, lisensi, pembatasan publikasi, dan perjanjian penggunaan data yang berlaku.
}

\newcommand{\PernyataanTanpaBiaya}{%
Apabila permintaan ini tidak dapat dipenuhi tanpa biaya, mohon kami diberi tahu terlebih dahulu. Kami tidak menyetujui pembuatan tabulasi atau pengadaan data berbayar tanpa persetujuan baru dari mahasiswa dan dosen pembimbing.
}

\newcommand{\AbstraksiPenelitian}{%
\textbf{Latar belakang.} Kegiatan ekonomi, perjalanan, layanan, dan perumahan di sekitar Jakarta berlangsung melampaui batas administrasi DKI Jakarta. Hubungan tersebut tidak otomatis menunjukkan bahwa pusat-pusat sekunder telah memperkuat fungsi lokalnya. Suatu wilayah dapat memperoleh akses ke pasar metropolitan, tetapi pada saat yang sama semakin bergantung pada Jakarta untuk pekerjaan, layanan, atau fungsi bernilai tinggi.

\textbf{Tujuan.} Penelitian ini mengkaji apakah integrasi dengan Jakarta berjalan bersama penguatan fungsi lokal, suatu pola yang konsisten dengan \textit{borrowed size}, atau bersama defisit fungsi dan ketergantungan yang asimetris, suatu pola yang konsisten dengan \textit{agglomeration shadow}. Penelitian tidak bertujuan mengusulkan perluasan wilayah DKI Jakarta, mengubah batas pemerintahan, atau membuktikan hubungan kausal hanya dari potret satu periode.

\textbf{Metode.} Kasus utama adalah sistem metropolitan Jabodetabek. Wilayah Statistik Metropolitan Indonesia digunakan sebagai sabuk pembanding. DKI Jakarta diperlakukan sebagai satu inti dengan batas administrasi tetap. Kecamatan menjadi unit utama jika resolusi data mendukung. Data pada tingkat kabupaten/kota, zona transportasi, raster, dan titik dipertahankan pada resolusi aslinya sebelum dilakukan agregasi yang dapat ditelusuri. Analisis membedakan data yang teramati, tabulasi agregat, hasil pemodelan, dan proksi.

\textbf{Cakupan wilayah dan waktu.} Baseline utama menggunakan tahun 2024. Data tahun lain digunakan sebagai konteks atau pemeriksaan perubahan, dengan tahun dan perbedaan unit ditampilkan secara terbuka. Wilayah mencakup DKI Jakarta, kabupaten/kota dalam kawasan utama Jabodetabek, pusat sekunder yang ditetapkan melalui prosedur yang terdokumentasi, serta sabuk pembanding yang diperlukan untuk menguji fungsi relatif.

\textbf{Jenis data dan variabel.} Penelitian memerlukan data penduduk, luas, kepadatan, integrasi dan perjalanan, fungsi layanan lokal, pekerjaan menurut lokasi tempat kerja dan sektor jika tersedia, aksesibilitas jaringan, serta karakteristik pusat sekunder. Podes digunakan untuk mengukur kehadiran dan akses layanan, bukan sebagai ukuran langsung produktivitas atau jumlah pekerjaan. Data komuter pada tingkat kabupaten/kota tidak akan dipaksa menjadi estimasi kecamatan. Data OD yang dimodelkan akan diberi label sintetis dan diuji terhadap kendala yang tersedia.

\textbf{Rencana hasil.} Hasil penelitian berupa profil integrasi dan fungsi lokal, tipologi kondisi relasional, peta gradien, uji sensitivitas bobot dan ambang, serta penjelajah skenario berbasis WebGL. Penjelajah tersebut menghitung ulang dan menampilkan hasil secara interaktif. Ia bukan sumber data waktu nyata, ramalan kebijakan, atau simulasi perubahan kewenangan pemerintahan.
}

\newcommand{\LampiranAbstraksi}{%
  \clearpage
  \section*{Lampiran 2. Abstraksi penggunaan data}
  \AbstraksiPenelitian
}

\newcommand{\TandaTanganDTAP}{%
  Hormat kami,\\[2.2cm]
  \parbox[t]{0.85\textwidth}{[NAMA PENANDATANGAN]}\\
  \parbox[t]{0.85\textwidth}{[JABATAN PENANDATANGAN SESUAI ALUR PERSURATAN DTAP]}
}

\newcommand{\TandaTanganBPS}{%
  Hormat kami,\\[2.2cm]
  \parbox[t]{0.85\textwidth}{[NAMA PENANDATANGAN]}\\
  \parbox[t]{0.85\textwidth}{[DEKAN FAKULTAS TEKNIK ATAU PEJABAT SETINGKAT YANG DITERIMA BPS]}
}


% Duplicate this draft only after the secondary centre and the receiving
% office have been selected and their official submission channel checked.
\newcommand{\KabupatenKota}{[KABUPATEN/KOTA]}
\newcommand{\AlamatInstansi}{[ALAMAT INSTANSI]}
\newcommand{\PusatSekunder}{[PUSAT SEKUNDER PRIORITAS]}
\newcommand{\WilayahTangkapan}{[WILAYAH TANGKAPAN]}
\newcommand{\PeriodeData}{2024 atau periode terdekat yang tersedia}

\begin{document}
\KopSurat

\begin{flushright}
Yogyakarta, [TANGGAL SURAT]
\end{flushright}

\begin{tabularx}{\textwidth}{@{}P{0.20\textwidth}Y@{}}
  Nomor & [DIISI OLEH UNIT PERSURATAN] \\
  Lampiran & Dua lampiran \\
  Hal & Permohonan tabulasi agregat pekerjaan menurut lokasi kerja
\end{tabularx}

Yth. Kepala Dinas Tenaga Kerja \KabupatenKota\\
\AlamatInstansi

Dengan hormat,

Sehubungan dengan penyusunan Tugas Akhir pada Program Sarjana Perencanaan Wilayah dan Kota, Departemen Teknik Arsitektur dan Perencanaan, Fakultas Teknik, Universitas Gadjah Mada, kami memohon dukungan data untuk mahasiswa berikut.

\IdentitasMahasiswa

Permohonan ini ditujukan untuk mendalami \PusatSekunder{} dan \WilayahTangkapan{}. Data pekerjaan menurut lokasi tempat kerja diperlukan agar kematangan ekonomi lokal tidak disimpulkan hanya dari jumlah penduduk, fasilitas, atau titik usaha pada peta.

Kami memohon tabulasi agregat yang telah tersedia dalam sistem Dinas Tenaga Kerja. Permintaan utama adalah jumlah pekerja menurut kecamatan lokasi tempat kerja dan kelompok lapangan usaha yang luas pada \PeriodeData. Jika tersedia, kami juga memohon indikator layanan ketenagakerjaan yang dapat membantu membedakan stok pekerjaan dari arus layanan. Data tidak perlu dibentuk secara khusus apabila hal tersebut menimbulkan biaya; bentuk agregat yang sudah tersedia tetap berguna sepanjang definisi dan cakupannya dijelaskan.

Permohonan ini tidak mencakup nama pekerja, NIK, alamat rumah, nomor kontak, atau identitas perusahaan yang dilindungi. Kami juga tidak menganggap data layanan ketenagakerjaan sebagai sensus seluruh pekerjaan sebelum cakupan pelaporan, periode, dan unit datanya diperiksa.

\PernyataanPenggunaanData

\PernyataanTanpaBiaya

Apabila sebagian data tidak dapat diberikan, kami bersedia menerima resolusi yang lebih kasar, tabulasi untuk \PusatSekunder{} dan \WilayahTangkapan{} saja, atau penjelasan tertulis mengenai unit terkecil yang dapat disajikan secara sah.

Demikian permohonan ini kami sampaikan. Atas perhatian dan bantuan Bapak/Ibu, kami mengucapkan terima kasih.

\vspace{8pt}
Tembusan:
\begin{enumerate}
  \item Dosen pembimbing;
  \item Ketua Program Studi Sarjana Perencanaan Wilayah dan Kota;
  \item Arsip.
\end{enumerate}

\TandaTanganDTAP

\clearpage
\section*{Lampiran 1. Rincian data yang dimohon}

\renewcommand{\arraystretch}{1.25}
\begin{tabularx}{\textwidth}{@{}P{0.12\textwidth}P{0.34\textwidth}Y@{}}
\toprule
Prioritas & Data yang dimohon & Unit, periode, dan kelengkapan \\
\midrule
P0 & Jumlah pekerja menurut kecamatan lokasi tempat kerja dan kelompok lapangan usaha luas & \PeriodeData; kode wilayah; definisi pekerja; sumber administrasi atau survei; dan cakupan usaha. \\
P0 & Jumlah lowongan, penempatan, atau peserta layanan ketenagakerjaan menurut kecamatan lokasi pekerjaan & Periode dan jenis layanan; bedakan stok dari arus; sertakan cakupan dan kemungkinan duplikasi. \\
P1 & Pekerja formal atau informal, atau tetap atau tidak tetap, menurut kecamatan jika definisi memadai & Kelompok luas dan aturan penyamaran sel kecil; tanpa data individu. \\
P1 & Jumlah pekerja menurut jenis kelamin atau kelompok umur luas jika dapat disajikan secara andal & Sertakan catatan keterwakilan, nilai hilang, dan aturan supresi. \\
P1 & Kamus data dan penjelasan kualitas & Definisi, periode pencatatan, unit, pembaruan, perubahan sistem, serta keterbatasan cakupan. \\
\bottomrule
\end{tabularx}

\vspace{8pt}
Jika data utama tidak tersedia, kami menerima salah satu bentuk pengganti berikut:
\begin{itemize}
  \item tabulasi untuk kecamatan yang membentuk \PusatSekunder{} dan \WilayahTangkapan{};
  \item agregasi kabupaten atau kota menurut sektor luas dengan penjelasan bahwa resolusi kecamatan tidak tersedia;
  \item indikator layanan ketenagakerjaan yang paling mendekati sisi tempat kerja, dengan batas cakupannya; atau
  \item penjelasan tertulis mengenai variabel yang tersedia tanpa pembuatan data berbayar.
\end{itemize}

Seluruh permintaan ditujukan pada data agregat atau data yang telah dianonimkan. Kami tidak meminta nama, alamat, nomor identitas, jejak perjalanan individu, atau pengenal usaha yang bersifat rahasia.

\LampiranAbstraksi

\end{document}
